\justifying
\NSFCBodyText

\subsubsection{项目特色}

本项目以新疆山地—盆地差异为研究条件，将模式大尺度状态、地形信息和多变量关系放在同一评价体系中考察。研究同时保留再分析内部配对实验与 CIESM 非配对应用两条证据线：前者判断条件重建是否有效，后者判断方法能否迁移到实际模式资料，并为数据处理集成提供依据。这一安排与一年期开放课题的资源相匹配，能够把方法验证和应用落实连接起来。

\subsubsection{创新点}

\textbf{（1）围绕地形条件建立变量联合增益的可归因检验。}
已有研究已开展多变量和时空生成降尺度。本项目拟进一步固定预测因子、偏差校正、模型容量和计算预算，对比地形条件下的独立与共享表征，研究地形分层如何改变变量联合的收益及负迁移。创新预期落在对局地条件与变量依赖关系的认识，以及这种认识对 CIESM 新疆区域重建的支持；若实验不能证明稳定增益，则据实给出适用边界。

\textbf{（2）分析大尺度一致性与局地有效细节的兼容边界。}
将校正后大尺度状态、输出投影和经验性地形约束分别处理，单独报告校正引起的信号变化、投影的数值残差及独立样本上的重建效果。拟揭示哪些约束能够限制不合理输出，哪些会压低极端或破坏变量共同变化，从而形成可选择、可检查的约束组合。该研究不把数学上强制满足的条件当作独立科学发现，而以受控对照中的收益和代价形成证据。
